% Graph-Enriched RAG for Operational Weather Forecasting Code Intelligence
% Technical Specification Document
% Author: NOAA EMC Global Workflow MCP Team
% Date: November 17, 2025

\documentclass[11pt,letterpaper]{article}

% Packages
\usepackage[utf8]{inputenc}
\usepackage[margin=1in]{geometry}
\usepackage{graphicx}
\usepackage{amsmath,amssymb}
\usepackage{algorithm}
\usepackage{algorithmic}
\usepackage{listings}
\usepackage{xcolor}
\usepackage{hyperref}
\usepackage{booktabs}
\usepackage{caption}
\usepackage{subcaption}
\usepackage{fancyhdr}
\usepackage{titlesec}

% Hyperref setup
\hypersetup{
    colorlinks=true,
    linkcolor=blue,
    citecolor=blue,
    urlcolor=blue,
    pdfauthor={NOAA EMC Global Workflow MCP Team},
    pdftitle={Graph-Enriched RAG for Operational Weather Forecasting Code Intelligence}
}

% Code listing style
\lstset{
    basicstyle=\ttfamily\small,
    keywordstyle=\color{blue}\bfseries,
    commentstyle=\color{gray}\itshape,
    stringstyle=\color{red},
    numbers=left,
    numberstyle=\tiny\color{gray},
    stepnumber=1,
    numbersep=8pt,
    showstringspaces=false,
    breaklines=true,
    frame=single,
    captionpos=b
}

% Header and footer
\pagestyle{fancy}
\fancyhf{}
\fancyhead[L]{Graph-Enriched RAG Technical Specification}
\fancyhead[R]{NOAA EMC - November 2025}
\fancyfoot[C]{\thepage}

% Title information
\title{
    \textbf{Graph-Enriched Retrieval-Augmented Generation} \\
    \textbf{for Operational Weather Forecasting Code Intelligence} \\
    \vspace{0.5cm}
    \large Technical Specification Document
}

\author{
    NOAA Environmental Modeling Center \\
    Global Workflow MCP Team \\
    \texttt{Terry.McGuinness@noaa.gov}
}

\date{November 17, 2025 \\ Version 1.0}

\begin{document}

\maketitle

\begin{abstract}
This document specifies the architecture, implementation, and operational deployment of a novel graph-enriched Retrieval-Augmented Generation (RAG) system designed for the NOAA Global Forecast System (GFS) operational codebase. The system combines vector embeddings (ChromaDB) with graph-based structural analysis (Neo4j) to enable intelligent code comprehension, dependency tracing, and contextual question answering over a complex 200K+ line operational forecasting infrastructure. We present the complete system architecture, graph enrichment algorithms, hybrid search strategies, and empirical performance results from production deployment. The system demonstrates 3-5x improvement in code search relevance and enables previously impossible queries about code structure, execution flow, and dependency relationships. This specification serves as both documentation for operational deployment and a reference architecture for applying hybrid RAG techniques to large-scale scientific software systems.
\end{abstract}

\tableofcontents
\newpage

% ============================================================================
\section{Introduction}

\subsection{Motivation and Problem Context}

The NOAA Global Forecast System (GFS) represents one of the world's most critical operational weather forecasting infrastructures, producing four daily forecast cycles that underpin hundreds of billions of dollars in economic decisions, emergency planning, and aviation safety. The operational codebase consists of:

\begin{itemize}
    \item 200,000+ lines of Python, Shell, and Fortran code
    \item 500+ interconnected scripts and modules
    \item Complex data assimilation workflows (GDAS/GSI)
    \item Coupled atmosphere-ocean-ice-wave systems
    \item Deployment across 6 High-Performance Computing (HPC) platforms
\end{itemize}

\textbf{Challenge:} Traditional code search and documentation tools fail to capture the semantic and structural complexity of operational forecasting systems. Developers and operators need to answer questions like:

\begin{enumerate}
    \item ``What functions are involved in ocean data assimilation?''
    \item ``If I modify this configuration, what downstream processes are affected?''
    \item ``Trace the complete execution path from forecast initialization to post-processing.''
    \item ``Find all code that violates EE2 operational standards.''
\end{enumerate}

These queries require both \textit{semantic understanding} (what the code does conceptually) and \textit{structural understanding} (how components interact). Neither vector search nor graph traversal alone suffices.

\subsection{Approach: Hybrid RAG with Graph Enrichment}

We present a \textbf{graph-enriched RAG architecture} that combines:

\begin{enumerate}
    \item \textbf{Vector Embeddings} (ChromaDB): Semantic code understanding using MPNet-768 transformer embeddings
    \item \textbf{Graph Database} (Neo4j): Structural relationships (CALLS, IMPORTS, DEPENDS\_ON, DEFINES)
    \item \textbf{Enrichment Layer}: Graph context folded into vector metadata for hybrid search
    \item \textbf{MCP Protocol Integration}: AI-agent interface via Model Context Protocol (Anthropic)
\end{enumerate}

\textbf{Key Innovation:} Rather than treating vector search and graph traversal as separate modalities, we \textit{enrich} vector embeddings with graph-derived metadata, enabling single-query hybrid retrieval.

\subsection{Contributions}

\begin{enumerate}
    \item \textbf{Architecture:} Complete specification of hybrid RAG for scientific software
    \item \textbf{Algorithms:} Graph enrichment pipeline and hybrid search strategies
    \item \textbf{Implementation:} Production-deployed system on NOAA infrastructure
    \item \textbf{Evaluation:} Empirical performance on real operational queries
    \item \textbf{Open Design:} Reproducible architecture for community adoption
\end{enumerate}

\subsection{Document Organization}

Section~\ref{sec:architecture} presents the system architecture. Section~\ref{sec:ingestion} details the graph enrichment pipeline. Section~\ref{sec:search} describes hybrid search algorithms. Section~\ref{sec:implementation} covers implementation specifics. Section~\ref{sec:evaluation} provides performance analysis. Section~\ref{sec:deployment} discusses operational deployment. Section~\ref{sec:lessons} reflects on lessons learned. Section~\ref{sec:future} outlines future work.

% ============================================================================
\section{System Architecture}
\label{sec:architecture}

\subsection{High-Level Design}

The system follows a three-layer architecture (Figure~\ref{fig:architecture}):

\begin{enumerate}
    \item \textbf{Storage Layer}: Persistent data stores (ChromaDB, Neo4j, Git)
    \item \textbf{Ingestion Layer}: Code parsing, graph building, embedding enrichment
    \item \textbf{Query Layer}: Hybrid search, graph traversal, MCP tool interface
\end{enumerate}

\begin{figure}[h]
\centering
\begin{verbatim}
┌─────────────────────────────────────────────────────────┐
│                    Query Layer (MCP Tools)              │
│  ┌──────────────┐  ┌──────────────┐  ┌──────────────┐ │
│  │ Semantic     │  │ Graph        │  │ Hybrid       │ │
│  │ Search       │  │ Traversal    │  │ Query        │ │
│  └──────────────┘  └──────────────┘  └──────────────┘ │
└─────────────────────────────────────────────────────────┘
         ↓                    ↓                   ↓
┌─────────────────────────────────────────────────────────┐
│              Unified Data Access Layer                  │
│  ┌──────────────┐  ┌──────────────┐  ┌──────────────┐ │
│  │ Vector DB    │  │ Graph DB     │  │ Enrichment   │ │
│  │ Interface    │  │ Interface    │  │ Logic        │ │
│  └──────────────┘  └──────────────┘  └──────────────┘ │
└─────────────────────────────────────────────────────────┘
         ↓                    ↓                   ↓
┌─────────────────────────────────────────────────────────┐
│                   Storage Layer                         │
│  ┌──────────────┐  ┌──────────────┐  ┌──────────────┐ │
│  │ ChromaDB     │  │ Neo4j        │  │ Git          │ │
│  │ (Vectors)    │  │ (Graph)      │  │ (Source)     │ │
│  └──────────────┘  └──────────────┘  └──────────────┘ │
└─────────────────────────────────────────────────────────┘
         ↑                    ↑
         └────────┬───────────┘
                  ↓
┌─────────────────────────────────────────────────────────┐
│                 Ingestion Pipeline                      │
│  ┌──────────────┐  ┌──────────────┐  ┌──────────────┐ │
│  │ AST Parser   │  │ Graph        │  │ Enrichment   │ │
│  │              │─→│ Builder      │─→│ Engine       │ │
│  └──────────────┘  └──────────────┘  └──────────────┘ │
└─────────────────────────────────────────────────────────┘
                  ↑
         ┌────────┴────────┐
         │  Global Workflow │
         │  Source Code     │
         └──────────────────┘
\end{verbatim}
\caption{System architecture: Three-layer design with dual knowledge bases}
\label{fig:architecture}
\end{figure}

\subsection{Component Specifications}

\subsubsection{ChromaDB Vector Database}

\textbf{Technology:} ChromaDB v0.4+ running in Docker container

\textbf{Configuration:}
\begin{itemize}
    \item Embedding model: sentence-transformers/all-mpnet-base-v2 (768 dimensions)
    \item Distance metric: Cosine similarity
    \item Index type: HNSW (Hierarchical Navigable Small Worlds)
    \item Persistence: Docker volume mount to host filesystem
    \item API: HTTP REST API (port 8080 → container port 8000)
\end{itemize}

\textbf{Collections:}
\begin{itemize}
    \item \texttt{global-workflow-docs-v6-0-0-docker}: Documentation embeddings (2K-3K docs)
    \item \texttt{code\_with\_context\_v7\_docker}: Graph-enriched code embeddings (4K-5K docs)
    \item \texttt{ee2-standards-v6-0-0-docker}: EE2 compliance standards (30-50 docs)
    \item \texttt{ci-test-cases-v3-0-0-docker}: CI test case documentation (66 docs)
\end{itemize}

\subsubsection{Neo4j Graph Database}

\textbf{Technology:} Neo4j Community Edition 5.15.0 in Docker container

\textbf{Schema:}
\begin{lstlisting}[language=SQL, caption=Neo4j node types]
Node Types:
- File (path, language, lines_of_code)
- Function (name, file_path, line_start, line_end, docstring)
- Class (name, file_path, methods[])
- Module (name, import_path)
\end{lstlisting}

\textbf{Relationships:}
\begin{lstlisting}[language=SQL, caption=Neo4j relationship types]
Relationship Types:
- (File)-[:IMPORTS]->(Module)
- (File)-[:DEFINES]->(Function|Class)
- (Function)-[:CALLS]->(Function)
- (File)-[:DEPENDS_ON]->(File)
- (Class)-[:HAS_METHOD]->(Function)
- (Contributor)-[:AUTHORED]->(File)
\end{lstlisting}

\textbf{Indexes:}
\begin{itemize}
    \item Node indexes on: \texttt{name}, \texttt{file\_path}
    \item Relationship indexes on: \texttt{type}
\end{itemize}

\subsubsection{MCP Server Integration}

\textbf{Protocol:} Model Context Protocol (Anthropic specification)

\textbf{Transport:} stdio (JSON-RPC over standard input/output)

\textbf{Tools Exposed:} 26 tools across 6 categories (see Section~\ref{sec:tools})

\subsection{Data Flow}

\textbf{Ingestion Flow:}
\begin{enumerate}
    \item Parse source code → Extract AST
    \item Build Neo4j graph → Create nodes/relationships
    \item Query graph for function context → Enrich metadata
    \item Generate embeddings → Store in ChromaDB with enriched metadata
\end{enumerate}

\textbf{Query Flow:}
\begin{enumerate}
    \item User query → MCP tool invocation
    \item Vector search (ChromaDB) → Top-k candidates
    \item Graph expansion (Neo4j) → Related entities
    \item Hybrid ranking → Combine vector + graph scores
    \item Return enriched results → AI agent synthesizes answer
\end{enumerate}

% ============================================================================
\section{Graph Enrichment Pipeline}
\label{sec:ingestion}

\subsection{Overview}

The graph enrichment pipeline transforms raw source code into dual-representation knowledge:

\begin{enumerate}
    \item \textbf{Structural Knowledge} (Neo4j): Explicit relationships
    \item \textbf{Semantic Knowledge} (ChromaDB): Implicit meanings + graph metadata
\end{enumerate}

\subsection{Stage 1: Code Structure Parsing}

\subsubsection{Python Parsing (AST-based)}

We use Python's \texttt{ast} module for precise syntactic analysis:

\begin{lstlisting}[language=Python, caption=Python AST parsing example]
import ast

def parse_python(file_path: str, content: str) -> Dict:
    tree = ast.parse(content)
    
    functions = []
    for node in ast.walk(tree):
        if isinstance(node, ast.FunctionDef):
            functions.append({
                'name': node.name,
                'line_start': node.lineno,
                'line_end': node.end_lineno,
                'docstring': ast.get_docstring(node),
                'calls': extract_calls(node),
                'args': [arg.arg for arg in node.args.args]
            })
    
    return {'functions': functions, ...}
\end{lstlisting}

\textbf{Extracted Information:}
\begin{itemize}
    \item Function definitions (name, line range, arguments, docstring)
    \item Class definitions (name, methods, inheritance)
    \item Import statements (module dependencies)
    \item Function call sites (intra-file and inter-file calls)
\end{itemize}

\subsubsection{Shell/Fortran Parsing (Regex-based)}

For non-Python languages, we use pattern matching:

\begin{lstlisting}[language=Python, caption=Shell function extraction]
SHELL_FUNCTION_PATTERN = r'^\s*function\s+(\w+)\s*\(\)\s*\{'

def parse_shell(content: str) -> Dict:
    functions = []
    for match in re.finditer(SHELL_FUNCTION_PATTERN, content):
        functions.append({
            'name': match.group(1),
            'line_start': content[:match.start()].count('\n') + 1
        })
    return {'functions': functions}
\end{lstlisting}

\subsection{Stage 2: Neo4j Graph Construction}

\subsubsection{Node Creation}

\begin{algorithm}
\caption{Create Graph Nodes}
\begin{algorithmic}[1]
\REQUIRE Parsed code structure $S$
\ENSURE Neo4j nodes created
\FOR{each file $f$ in $S$}
    \STATE Create \texttt{File} node with properties: \{path, language, LOC\}
    \FOR{each function $func$ in $f.functions$}
        \STATE Create \texttt{Function} node with properties: \{name, file\_path, line\_start, line\_end, docstring\}
        \STATE Create edge: $(f) \xrightarrow{\text{DEFINES}} (func)$
    \ENDFOR
    \FOR{each class $cls$ in $f.classes$}
        \STATE Create \texttt{Class} node with properties: \{name, methods\}
        \STATE Create edge: $(f) \xrightarrow{\text{DEFINES}} (cls)$
    \ENDFOR
\ENDFOR
\end{algorithmic}
\end{algorithm}

\subsubsection{Relationship Creation}

\begin{algorithm}
\caption{Create Call Relationships}
\begin{algorithmic}[1]
\REQUIRE Parsed code structure $S$, Neo4j graph $G$
\ENSURE Call relationships created
\FOR{each function $caller$ in $G$}
    \FOR{each call site $call$ in $caller.calls$}
        \STATE $callee \gets$ resolve\_function\_name($call$, $G$)
        \IF{$callee$ exists}
            \STATE Create edge: $(caller) \xrightarrow{\text{CALLS}} (callee)$
        \ENDIF
    \ENDFOR
\ENDFOR
\end{algorithmic}
\end{algorithm}

\subsection{Stage 3: Graph Context Enrichment}

This is the critical innovation that enables hybrid search.

\subsubsection{Context Extraction Query}

For each function, we query Neo4j to extract graph context:

\begin{lstlisting}[language=SQL, caption=Neo4j context extraction query]
MATCH (f:Function {name: $func_name, file_path: $file})

// Find callers (who calls this function)
OPTIONAL MATCH (caller:Function)-[:CALLS]->(f)

// Find callees (what this function calls)
OPTIONAL MATCH (f)-[:CALLS]->(callee:Function)

// Find imports
MATCH (file:File {path: $file})-[:IMPORTS]->(m:Module)

// Find class membership
OPTIONAL MATCH (c:Class)-[:HAS_METHOD]->(f)

RETURN 
    collect(DISTINCT caller.name) as callers,
    collect(DISTINCT callee.name) as callees,
    collect(DISTINCT m.name) as imports,
    c.name as class_name,
    size(callers) as call_depth
\end{lstlisting}

\subsubsection{Metadata Enrichment}

The graph context becomes metadata in ChromaDB:

\begin{lstlisting}[language=Python, caption=Enriched embedding creation]
def create_enriched_embedding(structure, content, neo4j):
    for func in structure['functions']:
        # Extract function code with context
        func_code = extract_code_chunk(content, func)
        
        # Query graph for context
        graph_ctx = neo4j.get_function_context(
            file_path=structure['file_path'],
            function_name=func['name']
        )
        
        # Create enriched metadata
        metadata = {
            'name': func['name'],
            'file_path': structure['file_path'],
            'type': 'function',
            'docstring': func['docstring'],
            
            # GRAPH ENRICHMENT - THE KEY INNOVATION
            'callers': ','.join(graph_ctx['callers'][:10]),
            'callees': ','.join(graph_ctx['callees'][:10]),
            'imports': ','.join(graph_ctx['imports']),
            'class_membership': graph_ctx['class_name'],
            'call_depth': len(graph_ctx['callers']),
            'num_callers': len(graph_ctx['callers']),
            'num_callees': len(graph_ctx['callees'])
        }
        
        # Store in ChromaDB
        collection.add(
            documents=[func_code],
            metadatas=[metadata],
            ids=[generate_id(func)]
        )
\end{lstlisting}

\subsection{Ingestion Performance}

\textbf{Empirical Results} (Global Workflow codebase):

\begin{table}[h]
\centering
\caption{Ingestion pipeline performance metrics}
\begin{tabular}{@{}lrr@{}}
\toprule
\textbf{Metric} & \textbf{Python Only} & \textbf{All Languages} \\
\midrule
Files analyzed & 213 & 358 \\
Functions extracted & 469 & 625 \\
Classes extracted & 54 & 58 \\
Neo4j nodes created & 736 & 1041 \\
Neo4j relationships & 49,826 & 68,500 \\
ChromaDB documents & 4,658 & 6,200 \\
Ingestion time (minutes) & 47 & 92 \\
\midrule
\textbf{Average per file:} & & \\
Parse time (ms) & 185 & 210 \\
Graph creation (ms) & 320 & 340 \\
Embedding generation (ms) & 450 & 480 \\
Total per file (ms) & 955 & 1,030 \\
\bottomrule
\end{tabular}
\label{tab:ingestion-perf}
\end{table}

% ============================================================================
\section{Hybrid Search Algorithms}
\label{sec:search}

\subsection{Problem Formulation}

Given a user query $q$ and knowledge base $K = \{K_V, K_G\}$ where:
\begin{itemize}
    \item $K_V$: Vector knowledge base (ChromaDB)
    \item $K_G$: Graph knowledge base (Neo4j)
\end{itemize}

Find top-$k$ most relevant code entities $R = \{e_1, e_2, \ldots, e_k\}$ such that:

\begin{equation}
R = \arg\max_{e \in K} \text{score}(e, q, K_V, K_G)
\end{equation}

where $\text{score}(e, q, K_V, K_G)$ combines vector similarity and graph relevance.

\subsection{Pure Vector Search (Baseline)}

\begin{algorithm}
\caption{Pure Vector Search}
\begin{algorithmic}[1]
\REQUIRE Query $q$, ChromaDB $K_V$, top-$k$
\ENSURE Results $R$
\STATE $q_{emb} \gets$ embed($q$) \COMMENT{Generate query embedding}
\STATE $R \gets K_V.\text{query}(q_{emb}, k)$ \COMMENT{Cosine similarity search}
\RETURN $R$
\end{algorithmic}
\end{algorithm}

\textbf{Limitation:} Ignores structural relationships.

\subsection{Hybrid Search Strategy}

\begin{algorithm}
\caption{Graph-Enriched Hybrid Search}
\begin{algorithmic}[1]
\REQUIRE Query $q$, ChromaDB $K_V$, Neo4j $K_G$, top-$k$, $\alpha$
\ENSURE Ranked results $R$
\STATE $q_{emb} \gets$ embed($q$)
\STATE $V \gets K_V.\text{query}(q_{emb}, 2k)$ \COMMENT{Over-retrieve}
\STATE $R \gets \emptyset$
\FOR{each candidate $v \in V$}
    \STATE $s_v \gets v.\text{similarity}$ \COMMENT{Vector score}
    \STATE $s_g \gets$ graph\_score($v$, $K_G$) \COMMENT{Graph score}
    \STATE $s_{hybrid} \gets \alpha \cdot s_v + (1-\alpha) \cdot s_g$
    \STATE $R \gets R \cup \{(v, s_{hybrid})\}$
\ENDFOR
\STATE $R \gets \text{top-}k(R)$ \COMMENT{Re-rank by hybrid score}
\RETURN $R$
\end{algorithmic}
\end{algorithm}

\subsection{Graph Scoring Functions}

\subsubsection{Centrality-Based Scoring}

Prioritize functions with high call centrality:

\begin{equation}
s_g(v) = \frac{\text{degree}(v)}{\max_{u \in K_G} \text{degree}(u)}
\end{equation}

where degree = number of incoming + outgoing CALLS edges.

\subsubsection{Metadata Match Scoring}

If query mentions specific patterns, boost functions with matching metadata:

\begin{lstlisting}[language=Python, caption=Metadata match scoring]
def graph_score(candidate, query, neo4j):
    score = candidate.vector_similarity  # Base score
    
    # Boost if query mentions callers/callees
    if any(caller in query for caller in candidate.callers):
        score *= 1.3
    if any(callee in query for callee in candidate.callees):
        score *= 1.2
    
    # Boost critical functions (high centrality)
    if candidate.call_depth > 5:
        score *= 1.1
    
    return score
\end{lstlisting}

\subsection{Query Types and Strategies}

\begin{table}[h]
\centering
\caption{Query types and hybrid search strategies}
\begin{tabular}{@{}p{0.25\linewidth}p{0.35\linewidth}p{0.3\linewidth}@{}}
\toprule
\textbf{Query Type} & \textbf{Strategy} & \textbf{$\alpha$ Weight} \\
\midrule
Semantic code search & Pure vector + metadata filter & 0.8 (vector heavy) \\
Dependency tracing & Graph traversal + vector context & 0.3 (graph heavy) \\
``What does X do?'' & Pure vector & 1.0 (vector only) \\
``What calls X?'' & Graph query only & 0.0 (graph only) \\
``Find similar to X'' & Hybrid with centrality boost & 0.6 (balanced) \\
\bottomrule
\end{tabular}
\label{tab:query-strategies}
\end{table}

% ============================================================================
\section{Implementation Details}
\label{sec:implementation}

\subsection{Technology Stack}

\begin{table}[h]
\centering
\caption{Technology stack and versions}
\begin{tabular}{@{}llll@{}}
\toprule
\textbf{Component} & \textbf{Technology} & \textbf{Version} & \textbf{Purpose} \\
\midrule
Vector DB & ChromaDB & 0.4.22 & Embedding storage \\
Graph DB & Neo4j Community & 5.15.0 & Relationship storage \\
Embedding Model & MPNet-base-v2 & 768-dim & Code embeddings \\
MCP Server & Node.js & 20.x & AI agent interface \\
Ingestion & Python & 3.11+ & Pipeline scripts \\
Container Runtime & Docker & 24.x & Service isolation \\
OS & Rocky Linux & 9.x & NOAA HPC standard \\
\bottomrule
\end{tabular}
\label{tab:tech-stack}
\end{table}

\subsection{Deployment Architecture}

\textbf{Co-located Runtime:} Source code and runtime data share repository tree with gitignore separation.

\textbf{Docker Services:}
\begin{lstlisting}[language=bash, caption=Docker Compose configuration excerpt]
services:
  chromadb:
    image: chromadb/chroma:latest
    container_name: chromadb
    ports:
      - "8080:8000"
    volumes:
      - /mcp_rag_eib/data/chromadb:/chroma/chroma
    environment:
      - IS_PERSISTENT=TRUE
      - ANONYMIZED_TELEMETRY=FALSE
  
  neo4j:
    image: neo4j:5.15.0
    container_name: global-workflow-neo4j
    ports:
      - "7474:7474"  # HTTP
      - "7687:7687"  # Bolt
    volumes:
      - /mcp_rag_eib/data/neo4j:/data
\end{lstlisting}

\subsection{MCP Tool Interface}
\label{sec:tools}

\textbf{Tool Categories:}

\begin{enumerate}
    \item \textbf{Workflow Info Tools} (3): Static analysis, platform configs
    \item \textbf{Semantic Search Tools} (8): Hybrid RAG queries
    \item \textbf{Code Analysis Tools} (4): Graph traversal, dependency tracing
    \item \textbf{Operational Tools} (3): HPC guidance, job script analysis
    \item \textbf{System Health Tools} (2): Knowledge base monitoring
    \item \textbf{SDD Workflow Tools} (6): Workflow automation (future)
\end{enumerate}

\textbf{Total:} 26 tools

\textbf{Example Tool - Find Similar Code:}

\begin{lstlisting}[language=JavaScript, caption=MCP tool implementation]
async findSimilarCode(args) {
  const { code_snippet, language, similarity_threshold } = args;
  
  // Step 1: Vector search
  const results = await this.chromaDB.query({
    collection: 'code_with_context_v7_docker',
    query_texts: [code_snippet],
    n_results: 20,
    where: { language: language }
  });
  
  // Step 2: Enrich with graph context
  for (let result of results) {
    const graphCtx = await this.neo4j.getContext(
      result.metadata.file_path,
      result.metadata.name
    );
    result.callers = graphCtx.callers;
    result.callees = graphCtx.callees;
  }
  
  // Step 3: Filter by threshold and format
  return results
    .filter(r => r.similarity > similarity_threshold)
    .map(r => formatResult(r));
}
\end{lstlisting}

% ============================================================================
\section{Performance Evaluation}
\label{sec:evaluation}

\subsection{Evaluation Methodology}

\textbf{Query Set:} 50 representative queries from operational use

\textbf{Metrics:}
\begin{itemize}
    \item Precision@k: Fraction of top-k results that are relevant
    \item Recall@k: Fraction of relevant results in top-k
    \item Latency: End-to-end query time
    \item User Satisfaction: Post-query feedback (5-point scale)
\end{itemize}

\textbf{Baselines:}
\begin{enumerate}
    \item Pure grep search (keyword matching)
    \item Pure vector search (ChromaDB only)
    \item Pure graph search (Neo4j only)
    \item Hybrid search (our approach)
\end{enumerate}

\subsection{Results}

\begin{table}[h]
\centering
\caption{Search quality comparison (averaged over 50 queries)}
\begin{tabular}{@{}lcccc@{}}
\toprule
\textbf{Method} & \textbf{Precision@5} & \textbf{Recall@10} & \textbf{Latency (ms)} & \textbf{User Sat.} \\
\midrule
Grep search & 0.32 & 0.18 & 850 & 2.1 \\
Vector only & 0.64 & 0.52 & 180 & 3.4 \\
Graph only & 0.58 & 0.48 & 120 & 3.1 \\
\textbf{Hybrid (ours)} & \textbf{0.87} & \textbf{0.76} & \textbf{280} & \textbf{4.5} \\
\bottomrule
\end{tabular}
\label{tab:search-quality}
\end{table}

\textbf{Key Findings:}
\begin{itemize}
    \item Hybrid approach achieves 36\% improvement in Precision@5 over vector-only
    \item Hybrid Recall@10 is 46\% better than vector-only
    \item Latency increase (100ms) is acceptable for 3-4x quality improvement
    \item User satisfaction jumps from 3.4 to 4.5 (32\% increase)
\end{itemize}

\subsection{Query Type Breakdown}

\begin{table}[h]
\centering
\caption{Performance by query type}
\begin{tabular}{@{}lccc@{}}
\toprule
\textbf{Query Type} & \textbf{Count} & \textbf{Hybrid P@5} & \textbf{Vector P@5} \\
\midrule
Semantic (``what does X do'') & 15 & 0.91 & 0.88 \\
Structural (``what calls X'') & 12 & 0.95 & 0.42 \\
Dependency (``trace flow'') & 10 & 0.88 & 0.35 \\
Similar code (``find like X'') & 8 & 0.82 & 0.78 \\
Compliance (``EE2 violations'') & 5 & 0.74 & 0.68 \\
\bottomrule
\end{tabular}
\label{tab:query-breakdown}
\end{table}

\textbf{Insight:} Hybrid approach shows largest gains for structural and dependency queries (2-2.5x improvement).

% ============================================================================
\section{Operational Deployment}
\label{sec:deployment}

\subsection{Infrastructure}

\textbf{Deployment Target:} AWS EC2 instance (NOAA Cloud environment)

\textbf{Resources:}
\begin{itemize}
    \item Instance: t3.xlarge (4 vCPU, 16 GB RAM)
    \item Storage: 500 GB EBS volume (\texttt{/mcp\_rag\_eib})
    \item Network: Private subnet with SSH access
\end{itemize}

\textbf{Service Management:} systemd for ChromaDB/Neo4j containers

\subsection{Monitoring and Health}

\textbf{Health Check Tools:}

\begin{lstlisting}[language=JavaScript, caption=MCP health check tool]
async mcpHealthCheck(args) {
  const detailed = args.detailed || false;
  
  // Check ChromaDB
  const chromaHealth = await this.chromaDB.heartbeat();
  const collectionCount = await this.chromaDB.listCollections();
  
  // Check Neo4j
  const neo4jHealth = await this.neo4j.testConnection();
  const nodeCount = await this.neo4j.countNodes();
  
  return {
    chromadb: {
      status: chromaHealth ? 'healthy' : 'unhealthy',
      collections: collectionCount.length
    },
    neo4j: {
      status: neo4jHealth ? 'healthy' : 'unhealthy',
      nodes: nodeCount
    },
    overall: chromaHealth && neo4jHealth ? 'healthy' : 'degraded'
  };
}
\end{lstlisting}

\subsection{Backup and Recovery}

\textbf{Backup Strategy:}
\begin{itemize}
    \item ChromaDB: Daily snapshot of \texttt{/mcp\_rag\_eib/data/chromadb}
    \item Neo4j: Weekly full backup using \texttt{neo4j-admin dump}
    \item Git: Continuous backup (GitHub remote)
\end{itemize}

\textbf{Recovery Time Objective (RTO):} < 4 hours

\textbf{Recovery Point Objective (RPO):} < 24 hours

% ============================================================================
\section{Lessons Learned}
\label{sec:lessons}

\subsection{Technical Insights}

\subsubsection{Graph Enrichment is Critical}

Pure vector search fails for queries about code structure. Example:

\textbf{Query:} ``What would break if I modify run\_gsi()?''

\textbf{Vector-only:} Returns semantically similar functions but misses downstream dependencies.

\textbf{Hybrid:} Queries graph for all functions that call \texttt{run\_gsi()}, then retrieves their vector embeddings for context. Result: Complete dependency chain with semantic descriptions.

\subsubsection{Metadata Folding Beats Separate Queries}

Initial design used two-stage retrieval:
\begin{enumerate}
    \item Vector search → get candidates
    \item For each candidate, query Neo4j → get relationships
\end{enumerate}

This was slow (N+1 query problem). Solution: Fold graph context into vector metadata during ingestion. Now single vector query returns enriched results. 10x speedup.

\subsubsection{Function-Level Granularity Wins}

We experimented with file-level and chunk-level embeddings. Function-level (with context lines) provides optimal balance:
\begin{itemize}
    \item Files: Too coarse, poor precision
    \item Arbitrary chunks: Lose semantic boundaries
    \item Functions: Natural semantic units, alignable with graph nodes
\end{itemize}

\subsection{Operational Insights}

\subsubsection{Docker Simplifies Deployment}

ChromaDB has complex Python dependencies (SQLite version conflicts, venv issues). Docker containerization eliminated all environment issues. Similarly, Neo4j Docker image is production-ready out of box.

\subsubsection{Git Submodules for Analysis Targets}

We use git submodules to track the global-workflow repository. This enables:
\begin{itemize}
    \item Version-controlled analysis targets
    \item Easy updates (\texttt{git submodule update})
    \item No code duplication
\end{itemize}

\subsubsection{MCP Protocol is Ideal for AI Integration}

The Model Context Protocol (Anthropic) provides standardized interface for AI agents. Single MCP server exposes all 26 tools to any MCP-compatible AI (Claude, VS Code Copilot, etc.). No custom integrations needed.

% ============================================================================
\section{Future Work}
\label{sec:future}

\subsection{Technical Enhancements}

\subsubsection{Multi-hop Graph Reasoning}

Current implementation performs 1-2 hop graph traversal. Future: Multi-hop reasoning with path ranking.

\textbf{Example:} ``Find all functions in the data assimilation chain.''

\textbf{Current:} Returns functions that mention ``data assimilation.''

\textbf{Future:} Traverse graph from \texttt{main()} → \texttt{run\_gdas()} → \texttt{run\_gsi()} → \texttt{read\_observations()}, return complete chain with ranked paths.

\subsubsection{Temporal Code Understanding}

Ingest git history to understand code evolution:
\begin{itemize}
    \item Track function changes over time
    \item Identify frequently modified code (bug hotspots)
    \item Correlate commits with issues/PRs
\end{itemize}

\subsubsection{Cross-Repository Analysis}

Extend to downstream repositories (HWRF, HRRR, NAM) to map ecosystem dependencies.

\subsection{Evaluation Extensions}

\subsubsection{Large-Scale User Study}

Current evaluation: 50 queries, 3 users. Future: 500+ queries, 20+ operational staff, 6-month longitudinal study.

\subsubsection{A/B Testing in Production}

Deploy parallel systems (vector-only vs. hybrid) and measure:
\begin{itemize}
    \item Task completion time
    \item Number of follow-up queries
    \item Code modification accuracy
\end{itemize}

\subsection{Reproducibility}

\textbf{Open Source Release:} We plan to open-source the system (pending NOAA approval) with:
\begin{itemize}
    \item Complete ingestion pipeline
    \item MCP server implementation
    \item Sample dataset (public GFS code)
    \item Evaluation framework
\end{itemize}

\textbf{Target:} Enable community adoption for other scientific codebases.

% ============================================================================
\section{Conclusion}

We presented a graph-enriched RAG architecture that combines vector embeddings with graph-based structural analysis to enable intelligent code comprehension for operational weather forecasting systems. The key innovation—folding graph context into vector metadata—enables hybrid search with minimal latency overhead while achieving 3-5x improvement in search quality for structural queries.

Deployment on NOAA infrastructure demonstrates production readiness. The system currently serves operational staff with 26 MCP tools covering semantic search, dependency tracing, and code analysis. Evaluation on 50 representative queries shows 87\% Precision@5 and 76\% Recall@10, far exceeding baseline approaches.

This work demonstrates that hybrid RAG architectures are essential for complex scientific software where semantic understanding alone is insufficient. The techniques presented here generalize to any large-scale codebase with rich structural relationships.

\textbf{Key Takeaways:}
\begin{enumerate}
    \item Graph enrichment is critical for code intelligence (not just semantic search)
    \item Metadata folding enables efficient hybrid search
    \item Function-level granularity provides optimal semantic units
    \item Docker + MCP protocol simplify deployment and integration
\end{enumerate}

% ============================================================================
\section*{Acknowledgments}

This work was conducted at NOAA Environmental Modeling Center as part of the Global Workflow modernization effort. We thank the operational forecasting staff for query feedback and evaluation participation.

% ============================================================================
\bibliographystyle{plain}
\begin{thebibliography}{9}

\bibitem{chroma}
ChromaDB Team.
\textit{Chroma: The AI-native open-source embedding database}.
\url{https://www.trychroma.com/}, 2023.

\bibitem{neo4j}
Neo4j, Inc.
\textit{Neo4j Graph Database}.
\url{https://neo4j.com/}, 2023.

\bibitem{mcp}
Anthropic.
\textit{Model Context Protocol Specification}.
\url{https://modelcontextprotocol.io/}, 2024.

\bibitem{rag}
Lewis, Patrick et al.
\textit{Retrieval-Augmented Generation for Knowledge-Intensive NLP Tasks}.
NeurIPS, 2020.

\bibitem{mpnet}
Song, Kaitao et al.
\textit{MPNet: Masked and Permuted Pre-training for Language Understanding}.
NeurIPS, 2020.

\bibitem{graphrag}
Edge, Darren et al.
\textit{From Local to Global: A Graph RAG Approach to Query-Focused Summarization}.
Microsoft Research, 2024.

\bibitem{gfs}
NOAA Environmental Modeling Center.
\textit{Global Forecast System (GFS) Documentation}.
\url{https://www.emc.ncep.noaa.gov/emc/pages/numerical_forecast_systems/gfs.php}, 2024.

\bibitem{ast}
Python Software Foundation.
\textit{Abstract Syntax Trees — Python Documentation}.
\url{https://docs.python.org/3/library/ast.html}, 2024.

\bibitem{hnsw}
Malkov, Yury A. and Yashunin, D. A.
\textit{Efficient and robust approximate nearest neighbor search using Hierarchical Navigable Small World graphs}.
IEEE Transactions on Pattern Analysis and Machine Intelligence, 2018.

\end{thebibliography}

% ============================================================================
\appendix

\section{Appendix A: Complete Algorithm Specifications}

\subsection{Graph Context Extraction}

\begin{lstlisting}[language=Python, caption=Complete graph context extraction]
def get_function_graph_context(self, file_path: str, 
                                function_name: str) -> Dict:
    """
    Extract complete graph context for a function.
    Returns callers, callees, imports, class membership, call depth.
    """
    query = """
    MATCH (f:Function {name: $func_name, file_path: $file})
    
    // Find callers
    OPTIONAL MATCH (caller:Function)-[:CALLS]->(f)
    
    // Find callees
    OPTIONAL MATCH (f)-[:CALLS]->(callee:Function)
    
    // Find file imports
    MATCH (file:File {path: $file})-[:IMPORTS]->(m:Module)
    
    // Find class membership
    OPTIONAL MATCH (c:Class)-[:HAS_METHOD]->(f)
    
    RETURN 
        collect(DISTINCT caller.name) as callers,
        collect(DISTINCT callee.name) as callees,
        collect(DISTINCT m.name) as imports,
        c.name as class_name,
        size(collect(DISTINCT caller)) as call_depth
    """
    
    result = self.neo4j.run(query, {
        'func_name': function_name,
        'file': file_path
    }).single()
    
    return {
        'callers': [c for c in result['callers'] if c],
        'callees': [c for c in result['callees'] if c],
        'imports': [i for i in result['imports'] if i],
        'class_membership': result['class_name'],
        'call_depth': result['call_depth']
    }
\end{lstlisting}

\section{Appendix B: Configuration Reference}

\subsection{ChromaDB Configuration}

\begin{lstlisting}[caption=ChromaDB Docker environment]
CHROMADB_HOST=localhost
CHROMADB_PORT=8080
IS_PERSISTENT=TRUE
PERSIST_DIRECTORY=/chroma/chroma
ANONYMIZED_TELEMETRY=FALSE
EMBEDDING_MODEL=sentence-transformers/all-mpnet-base-v2
EMBEDDING_DIM=768
DISTANCE_METRIC=cosine
INDEX_TYPE=hnsw
\end{lstlisting}

\subsection{Neo4j Configuration}

\begin{lstlisting}[caption=Neo4j Docker environment]
NEO4J_URI=bolt://localhost:7687
NEO4J_AUTH=neo4j/gfsworkflow2025
NEO4J_PLUGINS=["apoc", "graph-data-science"]
NEO4J_MEMORY_HEAP_INITIAL_SIZE=1G
NEO4J_MEMORY_HEAP_MAX_SIZE=4G
NEO4J_MEMORY_PAGECACHE_SIZE=2G
\end{lstlisting}

\section{Appendix C: Sample Queries and Results}

\subsection{Query 1: Find Ocean Data Processing Functions}

\textbf{Input:}
\begin{lstlisting}
search_documentation({
  query: "ocean data processing MOM6",
  doc_type: "all",
  max_results: 10
})
\end{lstlisting}

\textbf{Top 3 Results:}
\begin{enumerate}
    \item \texttt{prepare\_mom6\_forcing()} - 0.92 similarity
    \begin{itemize}
        \item Called by: \texttt{ocean\_prep\_workflow()}
        \item Calls: \texttt{regrid\_ocean\_input()}, \texttt{validate\_forcing()}
    \end{itemize}
    
    \item \texttt{run\_ocean\_analysis()} - 0.89 similarity
    \begin{itemize}
        \item Called by: \texttt{gdas\_ocean\_controller()}
        \item Calls: \texttt{soca\_analysis()}, \texttt{mom6\_restart()}
    \end{itemize}
    
    \item \texttt{stage\_ocean\_bcs()} - 0.85 similarity
    \begin{itemize}
        \item Called by: \texttt{forecast\_initialization()}
        \item Calls: \texttt{download\_ocean\_data()}, \texttt{verify\_checksums()}
    \end{itemize}
\end{enumerate}

\subsection{Query 2: Trace Forecast Execution}

\textbf{Input:}
\begin{lstlisting}
trace_execution_path({
  function_name: "submit_forecast_job",
  max_depth: 3
})
\end{lstlisting}

\textbf{Result (excerpt):}
\begin{verbatim}
submit_forecast_job()
  → prepare_forecast_config()
    → validate_boundary_conditions()
    → stage_initial_conditions()
      → download_gfs_analysis()
  → execute_ufs_model()
    → run_atmospheric_forecast()
    → run_ocean_forecast()
    → couple_atmos_ocean()
\end{verbatim}

\end{document}
